\chapter*{Appendix D: Story Scene Inventories}
\label{appendix:scenes}
\addcontentsline{toc}{chapter}{Appendix D: Story Scene Inventories}

The four tables in this appendix catalog each Story scene under the
sticky-graphic scrollytelling mechanic. Scene act is abbreviated
\textit{Bi} (Bing opener), \textit{B1}\ldots\textit{B$n$} (Bang
scenes), and \textit{Bo} (Bongo recap). The Tier column records
whether the pinned chart is a v5.5 catalog chart in story state
(Tier~1) or a bespoke simple story chart (Tier~2). See
Section~\ref{sec:solution:themes:stories} for the mechanic and
Table~\ref{tab:themes:chart_categories} for the two-tier sourcing rule.

\begin{table}[!htbp]
\centering
\caption{Story~1 scenes: \textit{The school commute, in numbers}.}
\label{tab:themes:story_school}
\footnotesize
\begin{tabularx}{\linewidth}{p{0.8cm}Xp{2.7cm}p{1.6cm}}
\hline
\textbf{Scene} & \textbf{Action title (scene opener)} & \textbf{Pinned chart (tier)} & \textbf{Transition} \\
\hline
Bing  & Thirty-six percent of HCMC bus rides are paid with a student-fare card.        & Bing hero frame                     & (entry)    \\
Bang  & Thirty-six of every hundred rides are student fares.                          & 2.3 weekly fare-class heatmap (Tier~1) & Highlight  \\
Bang  & Students compress their week into four hours.                                 & 2.3 weekly fare-class heatmap (Tier~1) & Filter     \\
Bang  & Saturday looks nothing like Tuesday.                                          & 2.3 weekly fare-class heatmap (Tier~1) & Filter     \\
Bang  & A handful of universities anchor most student demand.                         & 2.1 school constellation (Tier~1)     & Highlight  \\
Bang  & Three schools sit beyond walking distance of any stop.                        & 2.1 school constellation (Tier~1)     & Annotation \\
Bongo & Three schools without a stop, seven routes without a summer.                   & Bongo recap frame                   & (exit)     \\
\hline
\end{tabularx}
\end{table}

\begin{table}[!htbp]
\centering
\caption{Story~2 scenes: \textit{Where the city actually rides}.}
\label{tab:themes:story_spatial}
\footnotesize
\begin{tabularx}{\linewidth}{p{0.8cm}Xp{2.7cm}p{1.6cm}}
\hline
\textbf{Scene} & \textbf{Action title (scene opener)} & \textbf{Pinned chart (tier)} & \textbf{Transition} \\
\hline
Bing  & Five central stops carry seven percent of all boardings.                       & Bing hero frame                          & (entry)    \\
Bang  & The boarding map, at five-hundred-meter resolution.                            & 3.1 multi-layer hex map (Tier~1)         & Highlight  \\
Bang  & Sized by ridership, not by area, the center swells.                            & 3.1 multi-layer hex map (Tier~1)         & Filter     \\
Bang  & Five hex cells in the center dominate.                                         & 3.1 multi-layer hex map (Tier~1)         & Annotation \\
Bang  & Nineteen feeder routes connect to Metro Line~1.                                & 3.2 bus-to-metro Sankey (Tier~1)         & Highlight  \\
Bang  & Only eight arrive within five minutes of a train.                             & 3.2 bus-to-metro Sankey (Tier~1)         & Filter     \\
Bongo & The city you ride is not the city on the map.                                  & Bongo recap frame                        & (exit)     \\
\hline
\end{tabularx}
\end{table}

\begin{table}[!htbp]
\centering
\caption{Story~3 scenes: \textit{Who runs Saigon's buses}.}
\label{tab:themes:story_money}
\footnotesize
\begin{tabularx}{\linewidth}{p{0.8cm}Xp{2.7cm}p{1.6cm}}
\hline
\textbf{Scene} & \textbf{Action title (scene opener)} & \textbf{Pinned chart (tier)} & \textbf{Transition} \\
\hline
Bing  & \textdong{}72~billion in 90~days; two operators take 60\% of it.                & Bing hero frame                          & (entry)    \\
Bang  & Two companies move sixty percent of the network.                              & 4.2 operator constellation (Tier~1)      & Highlight  \\
Bang  & Size is not quality.                                                          & 4.2 operator constellation (Tier~1)      & Encode     \\
Bang  & The biggest operator runs 18\% of routes for 31\% of revenue.                  & 4.2 operator constellation (Tier~1)      & Annotation \\
Bang  & Twelve routes sit in the bleeding quadrant.                                   & 4.1 P\&L quadrant (Tier~1)               & Highlight  \\
Bang  & Free-fare student days lift ridership 12\% on weekdays.                        & free-fare lift slope (Tier~2)            & (entry)    \\
Bongo & Twelve bleeding routes; one slow digital transition; \textdong{}72B at stake. & Bongo recap frame                        & (exit)     \\
\hline
\end{tabularx}
\end{table}

\begin{table}[!htbp]
\centering
\caption{Story~4 scenes: \textit{Nine days without a network}.}
\label{tab:themes:story_tet}
\footnotesize
\begin{tabularx}{\linewidth}{p{0.8cm}Xp{2.7cm}p{1.6cm}}
\hline
\textbf{Scene} & \textbf{Action title (scene opener)} & \textbf{Pinned chart (tier)} & \textbf{Transition} \\
\hline
Bing  & Minus eighty-four percent: ridership on the trough day of T\^et.                & Bing hero frame                          & (entry)    \\
Bang  & Ridership held steady through January.                                         & 5.1 network pulse (Tier~1)               & (entry)    \\
Bang  & Then T\^et arrived, 14 February 2026.                                          & 5.1 network pulse (Tier~1)               & Annotation \\
Bang  & Minus eighty-four percent on the trough day, 16 February.                      & 5.1 network pulse (Tier~1)               & Annotation \\
Bang  & Recovery completes by 23 February.                                            & 5.1 network pulse (Tier~1)               & Annotation \\
Bang  & Some routes flatlined; others barely moved.                                    & resilience sparkline grid (Tier~2)       & (entry)    \\
Bongo & Nine days of flatline; nine routes need a contingency plan.                    & Bongo recap frame                        & (exit)     \\
\hline
\end{tabularx}
\par\smallskip
\footnotesize\textit{Note.} The T\^et percentages and dates were corrected
against the running warehouse; an earlier draft placed the window at
January~28--February~12 with $-67\%$. The measured window runs
February~14--22 with a $-57\%$ average and $-84\%$ on February~16.
\end{table}
